\section*{Bab \ref{ch:g10:lines} --- Persamaan Garis dan Sistem Persamaan Linear}

\begin{solution}{exo:g10:lines:1}
$y = 2x - 3$: \omterm{def:g10:reffunc:affine}{gradien} $2$, titik potong $-3$.
$y = -\frac13 x + 1$: \omterm{def:g10:reffunc:affine}{gradien} $-\frac13$, titik potong $1$.
$y = 4$: \omterm{def:g10:reffunc:affine}{gradien} $0$, titik potong $4$ (garis mendatar).
$x = -2$: garis tegak, tanpa \omterm{def:g10:reffunc:affine}{gradien} dan tanpa \omterm{thm:g10:lines:reduced}{persamaan tereduksi}.
\end{solution}

\begin{solution}{exo:g10:lines:2}
$3 \times 2 + 1 = 7$: ya, $A$ terletak pada garis itu.
$3 \times (-1) + 1 = -2 \neq -1$: $B$ tidak.
$3 \times 4 + 1 = 13$: ya, $C$ terletak pada garis itu.
\end{solution}

\begin{solution}{exo:g10:lines:3}
(a) $m = \frac{8 - 2}{3 - 0} = 2$ dan $p = y_A = 2$ (titik $A$ terletak pada
sumbu $y$): $y = 2x + 2$.

(b) $m = \frac{-4 - 5}{2 - (-1)} = \frac{-9}{3} = -3$; lalu
$5 = -3 \times (-1) + p$ memberi $p = 2$: $y = -3x + 2$.

(c) $x_A = x_B = 4$: garis tegak $x = 4$.
\end{solution}

\begin{solution}{exo:g10:lines:4}
$y = \frac{6x+2}{2} = 3x + 1$. Garis yang bergradien $3$ adalah $y = 3x - 1$,
$y = 3x + 5$ dan $y = 3x + 1$: ketiganya saling sejajar;
tidak ada dua di antaranya yang berimpit (titik potongnya berbeda), dan tidak
satu pun sejajar dengan $y = -3x + 2$.
\end{solution}

\begin{solution}{exo:g10:lines:5}
Sistem pertama: substitusikan $y = 2x - 1$ ke $3x + y = 9$:
$3x + 2x - 1 = 9$, jadi $5x = 10$, $x = 2$, lalu $y = 3$. Penyelesaiannya:
$(2, 3)$.

Sistem kedua: dari $x - y = 4$, diperoleh $x = y + 4$. Lalu
$2(y + 4) + 3y = 3$, jadi $5y = -5$, $y = -1$, lalu $x = 3$. Penyelesaiannya:
$(3, -1)$.
\end{solution}

\begin{solution}{exo:g10:lines:6}
Sistem pertama: menjumlahkan kedua \omterm{def:g10:algebra:equation}{persamaan} melenyapkan $y$: $2x = 14$,
$x = 7$; lalu $7 + y = 10$ memberi $y = 3$. Penyelesaiannya: $(7, 3)$.

Sistem kedua: kalikan \omterm{def:g10:algebra:equation}{persamaan} pertama dengan $5$ dan yang kedua dengan $-2$:
\[
\begin{cases}
15x + 10y = 60 \\
-4x - 10y = -16
\end{cases}
\]
Setelah dijumlahkan: $11x = 44$, jadi $x = 4$; lalu $3 \times 4 + 2y = 12$
memberi $y = 0$. Penyelesaiannya: $(4, 0)$.
\end{solution}

\begin{solution}{exo:g10:lines:7}
Hitunglah $ab' - a'b$ pada setiap kasus.

Pertama: $2 \times 2 - (-4) \times (-1) = 4 - 4 = 0$, dan \omterm{def:g10:algebra:equation}{persamaan}
keduanya sama dengan $-2$ kali yang pertama: dua kali garis yang sama, jadi
tak hingga banyak penyelesaian.

Kedua: $ab' - a'b = 0$ lagi, tetapi ruas kanannya tidak berbanding
$-2$ ($1 \neq -6$): dua garis sejajar yang berbeda, jadi tanpa penyelesaian.

Ketiga: $2 \times 1 - 1 \times (-1) = 3 \neq 0$: tepat satu penyelesaian.
\end{solution}

\begin{solution}{exo:g10:lines:8}
Misalkan $x$ banyaknya karcis anak-anak dan $y$ banyaknya karcis
dewasa:
\[
\begin{cases}
x + y = 200 \\
5x + 9y = 1352 .
\end{cases}
\]
Dari \omterm{def:g10:algebra:equation}{persamaan} pertama $x = 200 - y$; setelah disubstitusikan,
$5(200 - y) + 9y = 1352$, jadi $1000 + 4y = 1352$, $y = 88$, lalu
$x = 112$. Jadi $112$ karcis anak-anak dan $88$ karcis dewasa. Periksa:
$5 \times 112 + 9 \times 88 = 560 + 792 = 1352$.
\end{solution}

\begin{solution}{exo:g10:lines:9}
\emph{1.} Sejajar berarti \omterm{def:g10:reffunc:affine}{gradiennya} sama, yaitu $2$: $y = 2x + p$ dengan
$4 = 2 \times 1 + p$, jadi $p = 2$: $d'\colon y = 2x + 2$.

\emph{2.} Pada sumbu $x$ berlaku $y = 0$: $2x + 2 = 0$ memberi $x = -1$.
Titik potongnya adalah $(-1, 0)$.
\end{solution}

\begin{solution}{exo:g10:lines:10}
\emph{1.} \omterm{prop:g10:coordgeom:midpoint}{Titik tengah} $[BC]$: $A' = \left(\frac{6+2}{2},
\frac{2+4}{2}\right) = (4, 3)$. Garis berat dari $A(0,0)$ bergradien
$\frac{3 - 0}{4 - 0} = \frac34$: \omterm{def:g10:algebra:equation}{persamaannya} $y = \frac34 x$.

\omterm{prop:g10:coordgeom:midpoint}{Titik tengah} $[AC]$: $B' = (1, 2)$. Garis berat dari $B(6, 2)$ bergradien
$\frac{2 - 2}{1 - 6} = 0$: garis itu mendatar, $y = 2$.

\emph{2.} Titik potongnya: $\frac34 x = 2$ memberi $x = \frac83$, jadi
$G\left(\frac83, 2\right)$. Garis berat ketiga menghubungkan $C(2,4)$ dengan
\omterm{prop:g10:coordgeom:midpoint}{titik tengah} $[AB]$, yaitu $C' = (3, 1)$; \omterm{def:g10:reffunc:affine}{gradiennya}
$\frac{1 - 4}{3 - 2} = -3$, \omterm{def:g10:algebra:equation}{persamaannya} $y = -3x + 10$. Di $x = \frac83$:
$y = -8 + 10 = 2$: jadi $G$ terletak padanya. Dan memang
$G = \left(\frac{0 + 6 + 2}{3}, \frac{0 + 2 + 4}{3}\right)$: titik
beratnya merata-ratakan \omterm{def:g10:coordgeom:system}{koordinat} ketiga titik sudutnya.
\end{solution}

\begin{solution}{pb:g10:lines:1}
\textbf{1.} \omterm{def:g10:reffunc:affine}{Gradiennya} $\frac{8 - 2}{3 - 1} = 3$; melalui $(1, 2)$:
$y = 3x - 1$.

\textbf{2.} $y = 3x - 1$ dan $y = 3x + 4$: \omterm{def:g10:reffunc:affine}{gradiennya} sama, jadi sejajar
(dan berbeda). Memotongkan $y = 3x - 1$ dengan $y = -x + 7$:
$3x - 1 = -x + 7$ memberi $x = 2$, $y = 5$: titik $(2, 5)$.

\textbf{3.} Kedua titiknya berbagi $x = 2$: garisnya tegak,
\omterm{def:g10:algebra:equation}{persamaannya} $x = 2$. Bentuk tereduksi $y = mx + p$ menjawab ``satu $y$
untuk setiap $x$'', dan itu ditolak oleh garis tegak: \omterm{def:g10:reffunc:affine}{gradiennya} akan
menjadi takhingga.

\textbf{4.} $3 \times 4 - 1 = 11$: ya, $(4, 11)$ terletak pada
garis itu. $3 \times (-1) - 1 = -4 \neq -5$: tidak.

\textbf{5.} $y = -2x + 3$; memotong kedua sumbu di $(0, 3)$ dan
$\left(\frac32, 0\right)$. Luas segitiganya:
$\frac12 \times \frac32 \times 3 = \frac94 = 2.25$.

\textbf{6.} Setelah disubstitusikan: $3x + 2x - 3 = 7$, $x = 2$,
$y = 1$: penyelesaiannya $(2, 1)$.

\textbf{7.} Setelah dijumlahkan: $7x = 7$, $x = 1$; lalu $3y = 6$, $y = 2$:
penyelesaiannya $(1, 2)$.

\textbf{8.} Sistem pertama: \omterm{def:g10:algebra:equation}{persamaan} keduanya dua kali yang
pertama: satu garis yang terhitung dua kali --- tak hingga banyak
penyelesaian (semua titik pada $2x - y = 3$). Sistem kedua: ruas kirinya
sebanding, ruas kanannya tidak ($2 \times 3 = 6 \neq 1$): dua
garis sejajar --- tanpa penyelesaian. Trikotominya: \omterm{def:g10:reffunc:affine}{gradien} berbeda
$\to$ satu titik potong; \omterm{def:g10:reffunc:affine}{gradien} sama, titik potong sumbu berbeda $\to$
sejajar, tanpa penyelesaian; semuanya sama $\to$ garis yang sama, tak hingga
banyak penyelesaian.

\textbf{9.} \omterm{def:g10:vectors:vector}{Vektor} arah kedua garisnya adalah $(-b, a)$ dan
$(-d, c)$ (atau: \omterm{def:g10:reffunc:affine}{gradiennya} $-\frac ab$ dan $-\frac cd$), yang \omterm{def:g10:vectors:collinear}{kolinear}
tepat ketika $ad - bc = 0$ --- determinan
\cref{pb:g10:vectors:1} dalam pekerjaan baru. Nilainya: pertanyaan 7:
$2 \times (-3) - 3 \times 5 = -21 \neq 0$: penyelesaiannya tunggal.
Pertanyaan 8: $2 \times (-2) - (-1) \times 4 = 0$ untuk keduanya:
sejajar atau berimpit, dan konstantanyalah yang memisahkan kedua
kasus itu.

\textbf{10.} $ad - bc = 6 - 2m$: bernilai nol untuk $m = 3$. Di sana
$x + 3y = 3$ berhadapan dengan $2x + 6y = 7$: melipatduakan yang pertama
memberi $2x + 6y = 6 \neq 7$: sejajar, jadi \emph{tanpa penyelesaian}. Untuk
$m \neq 3$, eliminasi memberi penyelesaian tunggal
$y = \frac{1}{6 - 2m}$, lalu $x = 3 - \frac{m}{6 - 2m}$.

\textbf{11.} $p + r = 35$, $2p + 4r = 94$: setelah dibagi dua,
$p + 2r = 47$, jadi $r = 12$ dan $p = 23$: dua puluh tiga burung pegar,
dua belas kelinci --- jawaban kitab \emph{Sembilan Bab} sendiri.

\textbf{12.} $x + y = 30$ dan $0.6x + 0.2y = 0.35 \times 30 =
10.5$: dengan mengurangkan $0.2 \times$ yang pertama, $0.4x = 4.5$:
$x = 11.25$ L sirop, $y = 18.75$ L minuman.

\textbf{13.} Angkanya $a$, $b$: $a + b = 12$ dan
$(10a + b) - (10b + a) = 9(a - b) = 18$, jadi $a - b = 2$:
$a = 7$, $b = 5$: bilangannya $75$.

\textbf{14.} Setelah dijumlahkan: $5c + 5r = 18$, jadi $c + r = 3.60$ ---
yaitu satu kopi ditambah satu croissant. Setelah dikurangkan: $c - r = 0.20$.
Jadi $c = 1.90$ dan $r = 1.70$ euro: jalan pintas yang setangkup itu
menyelesaikannya tanpa substitusi apa pun.

\textbf{15.} Melipatduakan pengakuan pertama memberi $4$ apel $+ 2$
pir $= 10$ euro, padahal lapak itu menarik $9$: sistemnya
taktaat-asas --- garis sejajar, himpunan penyelesaiannya kosong. Putusannya:
kedua pengakuan harga itu tidak mungkin benar bersamaan; entah ada potongan
harga yang disembunyikan, entah aritmetikanya.

\textbf{16.} Di titik asal: $0 > 3 \times 0 - 1 = -1$: benar,
jadi $(0,0)$ terletak di daerah $y > 3x - 1$ (di atas garisnya).
Ujilah sebarang titik dengan memasukkannya: di atas atau di bawah menurut
benar tidaknya pertidaksamaan itu.

\textbf{17.} Titik sudutnya: $(0, 0)$; $(4, 0)$ (sumbu dan
$x + y = 4$); $(0, 1)$ (sumbu dan $y = 2x + 1$); serta perpotongan
$x + y = 4$ dengan $y = 2x + 1$: $x = 1$, $y = 3$: yaitu $(1, 3)$.

\textbf{18.} $P(0,0) = 0$; $P(4, 0) = 12$; $P(0, 1) = 2$;
$P(1, 3) = 9$. \omterm{def:g10:functions:extrema}{Maksimumnya} di titik sudut $(4, 0)$: rencana
$x = 4$, $y = 0$ memberi $12$ --- titik sudutlah yang menentukan, seperti
diperlihatkan garis sejajar $3x + 2y = c$ yang menyapu daerahnya.

\textbf{19.} $1.5x + 3 = 2x$ memberi $x = 6$, $y = 12$: pada enam
kilometer kedua taksi menarik dua belas euro --- titik impas pada
soal akhir pekan tentang dua termometer di jilid sebelumnya, kini menjadi
penyelesaian tunggal sebuah sistem.

\textbf{20.} Penyelesaian sebuah \omterm{def:g10:algebra:equation}{persamaan} linear menggambar sebuah garis;
sebuah sistem menumpukkan dua garis, dan $ad - bc$ memberitahukan sekilas
apakah keduanya berpotongan sekali (bila taknol) atau jatuh ke kasus
sejajar atau berimpit (bila nol). Tumpuklah pertidaksamaan sebagai gantinya
dan penyelesaiannya membentuk daerah segi banyak yang titik sudutnya
mengusung nilai optimum linear mana pun. Aljabar linear memperumum
determinan itu ke sebarang banyak variabel; pemrograman linear
memperindustrikan titik sudutnya.

\end{solution}
